% docs/texte/10_engineering/data_formats.tex -- chapter "Data Formats" (Interferosim)
% Version: 0.4.0 (2026-10-31) -- VisData in memory
% 0.3.0 (2026-10-31) -- receiver keys of the device profile, base
% 0.2.0 (2026-10-12) -- arrays: site/antenna/receiver references,
%   stations as offsets east/north/up; horn design_freq_hz
%   0.1.0: first version
% ================================================================
\label{ch:data-formats}
\index{data formats}

\begin{center}
\begin{tabular}{@{}lll@{}}
\toprule
Content & Format & Location \\
\midrule
Sky models & JSON (AAS6 \texttt{sim\_sky.json}) & \texttt{configs/skies/} \\
Antennas & TOML & \texttt{configs/antennas/} \\
Receivers & TOML & \texttt{configs/receivers/} \\
Sites and stations & TOML & \texttt{configs/sites/} \\
Visibilities & FITS binary table & output of the simulation \\
Maps, spectra, time series & FITS (from \texttt{aas6sim}) & output \\
\bottomrule
\end{tabular}
\end{center}

\section{Sky models}

See chapter~\ref{ch:sky-models}. JSON with \texttt{center}
(\texttt{label}, \texttt{ra\_deg}, \texttt{dec\_deg}) and
\texttt{components} (\texttt{name}, \texttt{dx\_deg}, \texttt{dy\_deg},
\texttt{flux\_jy}, \texttt{kind} and the parameters of the kind).

\section{Configuration files (TOML)}

\paragraph*{Antennas} (\texttt{configs/antennas/*.toml}), common keys
\texttt{name}, \texttt{type}; further keys depend on the type:
\begin{center}
\begin{tabular}{@{}ll@{}}
\toprule
\texttt{type} & Keys \\
\midrule
\texttt{horn} & \texttt{gain\_dbi}, optional \texttt{design\_freq\_hz} (fixed aperture) \\
\texttt{dish} & \texttt{diameter\_m}, \texttt{aperture\_efficiency}, \texttt{beam\_factor} \\
\texttt{yagi} & \texttt{n\_elements}, \texttt{boom\_length\_m}, \texttt{gain\_dbi},
                \texttt{design\_freq\_hz}, \texttt{polarisation} \\
\bottomrule
\end{tabular}
\end{center}

\paragraph*{Receivers} (\texttt{configs/receivers/*.toml}):
\texttt{name}, \texttt{rx\_channels}, \texttt{sample\_rate\_sps},
\texttt{bandwidth\_hz}, \texttt{bits}, \texttt{t\_sys\_k}, \texttt{agc},
\texttt{dc\_offset\_correction}, \texttt{iq\_balance\_correction},
\texttt{clock\_source}; since 0.3.2 also \texttt{center\_freq\_hz},
\texttt{rx\_gain\_db}, \texttt{sample\_format}, \texttt{time\_source},
\texttt{external\_lna}, the device data of chapter~\ref{ch:receiver}
(\texttt{freq\_min\_hz} \dots\ \texttt{fpga}) and \texttt{base} (name of a
profile in the same directory whose keys are read first).

\paragraph*{Sites and stations} (\texttt{configs/sites/*.toml}): a single
site has \texttt{name}, \texttt{latitude\_deg}, \texttt{longitude\_deg},
\texttt{altitude\_m}; a group of stations has \texttt{reference} and one
table \texttt{[stations.<id>]} per station with \texttt{latitude\_deg},
\texttt{longitude\_deg}, \texttt{height\_m} (WGS84). The two stations of
the AAS6 simulation are in \texttt{aas6\_sim\_stations.toml}
(from \texttt{sb\_iq\_sim\_a/b.json}).
Instead of the geodetic position a station may be given as offset from
the site in metres (\texttt{east\_m}, \texttt{north\_m}, \texttt{up\_m}),
e.\,g.\ \texttt{two\_horns\_10m.toml}. A group of stations may name
further configuration files by their name without \texttt{.toml}:
\texttt{site} (in \texttt{configs/sites/}), \texttt{antenna} (in
\texttt{configs/antennas/}, also per station) and \texttt{receiver} (in
\texttt{configs/receivers/}); \texttt{Array.from\_toml()} reads them
(chapter~\ref{ch:array}).

\section{Visibilities (FITS)}

Format of \texttt{aas6sim} (\texttt{visibility.py}): primary HDU without
data plus a binary table \texttt{EXTNAME = VISIBILITY}, one row per
integration.
\begin{center}
\begin{tabular}{@{}lll@{}}
\toprule
Column & Unit & Content \\
\midrule
\texttt{TIME} & d & MJD (UTC) of the middle of the interval \\
\texttt{HA} & deg & hour angle \\
\texttt{UU}, \texttt{VV}, \texttt{WW} & wavelengths & baseline components \\
\texttt{VISREAL}, \texttt{VISIMAG} & Jy & real and imaginary part \\
\texttt{WEIGHT} & Jy$^{-2}$ & weight $1/\sigma^2$ \\
\bottomrule
\end{tabular}
\end{center}
Header keywords: \texttt{ORIGIN}, \texttt{DATE}, \texttt{SIMULATD}
(\texttt{T} for simulated data -- a simulated file must remain
recognisable after any number of processing steps), \texttt{TELESCOP},
\texttt{INSTRUME}, \texttt{OBJECT}, \texttt{DATE-OBS}, \texttt{UVUNIT},
\texttt{ANT1}, \texttt{ANT2}, \texttt{BASELINE} and the observation keys
\texttt{RESTFRQ}, \texttt{BANDWID}, \texttt{TSYS}, \texttt{EXPOSURE},
\texttt{OBSRA}, \texttt{OBSDEC}, \texttt{SITELAT}, \texttt{SITELONG},
\texttt{SITEELEV}.

\paragraph*{Extension for the frequency axis.} With several channels
(F-SKY-5, F-DAT-1) the columns \texttt{VISREAL}, \texttt{VISIMAG} and
\texttt{WEIGHT} become vectors of length $n_{ch}$ per row, and a
second table \texttt{EXTNAME = FREQUENCY} holds the channel
frequencies in Hz. A file with one channel remains identical to the
\texttt{aas6sim} format.

\paragraph*{Data class in memory.} \texttt{interferosim.visdata.VisData}
(since 0.5.0) holds the same columns as arrays -- \texttt{time\_mjd},
\texttt{ha\_deg}, \texttt{u}, \texttt{v}, \texttt{w}, \texttt{vis} (rows
$\times$ channels, complex), \texttt{weight} -- plus a baseline index per row
with \texttt{baseline\_names} (\texttt{ANT1}/\texttt{ANT2} of the file), the
channel frequencies \texttt{freq\_hz} and the header keywords as a
dictionary. Until the FITS output of the package \texttt{io} exists,
\texttt{save\_npz}/\texttt{load\_npz} store it losslessly.

\section{Further FITS files}

Taken over unchanged from \texttt{aas6sim}: maps (2D image with WCS,
\texttt{skymap.py}), cross-spectra cubes (\texttt{spectra.py}),
measurement time series (\texttt{timeseries.py}) and an UVFITS export
(\texttt{uvfits.py}, random-groups without AN/FQ tables; acceptance by
AIPS, CASA or difmap not tested).
